% Part: lambda-calculus
% Chapter: lambda-definability
% Section: truth-values

\documentclass[../../../include/open-logic-section]{subfiles}

\begin{document}

\olfileid{lam}{ldf}{tvr}
\olsection{قيم الصدق والعلاقات}

يمكننا ترميز قيم الصدق في حساب لامبدا الخالص كما يأتي:
\begin{align*}
  \fn{true} & \ident \lambd[x][\lambd[y][x]]\\
  \fn{false} & \ident \lambd[x][\lambd[y][y]]
\end{align*}

تمثل قيم الصدق بوصفها \emph{دوال اختيار}، أي دوال تقبل وسيطين وتعيد
أحدهما. فتختار قيمة الصدق~$\fn{true}$ وسيطها الأول، وتختار~$\fn{false}$
وسيطها الثاني. فمثلًا، يختزل~$\fn{true}\, M N$ دائمًا إلى~$M$، في حين
يختزل~$\fn{false}\, M N$ دائمًا إلى~$N$.

\begin{defn}
نقول إن العلاقة~$R \subseteq \Nat^k$ !!{lambda definable} إذا وجد
حد~$R$ بحيث
\begin{align*}
  R\, \num{n_1} \dots \num{n_k} & \bred \fn{true}
  \intertext{متى تحقق~$R(n_1, \dots, n_k)$، و}
  R\, \num{n_1} \dots \num{n_k} & \bred \fn{false}
\end{align*}
في غير ذلك.
\end{defn}

فمثلًا، العلاقة~$\fn{IsZero} = \{0\}$ التي تتحقق في~$0$ وحده
!!{lambda definable} بالحد
\[
  \fn{IsZero} \ident \lambd[n][n (\lambd[x][\fn{false}])\, \fn{true}].
\]
كيف يعمل؟ بما أن أعداد تشيرش تعرّف بوصفها مكررات (أي دوال تطبق وسيطها
الأول~$n$ مرة في الثاني)، نجعل القيمة الابتدائية~$\fn{true}$، ونعيد
$\fn{false}$ في كل خطوة من التكرار بصرف النظر عن نتيجة الخطوة السابقة.
ستطبق هذه الخطوة في القيمة الابتدائية~$n$ مرة، وستكون النتيجة
$\fn{true}$ إذا وفقط إذا لم تطبق الخطوة أصلًا، أي عندما~$n = 0$.

وبناء على هذا التمثيل لقيم الصدق، يمكننا تعريف بعض دوال الصدق الأخرى.
وفيما يأتي دالتان تمثلان النفي والاقتران:
\begin{align*}
  \fn{Not} & \ident \lambd[x][x\, \fn{false}\, \fn{true}]\\
  \fn{And} & \ident \lambd[x][\lambd[y][xy \,\fn{false}]]
\end{align*}
تقبل الدالة ``$\fn{Not}$'' وسيطًا واحدًا، وتعيد~$\fn{true}$ إذا كان
الوسيط~$\fn{false}$، وتعيد~$\fn{false}$ إذا كان الوسيط~$\fn{true}$.
أما الدالة ``$\fn{And}$'' فتقبل قيمتي صدق وسيطين، وينبغي أن تعيد
$\fn{true}$ إذا وفقط إذا كان كلا الوسيطين~$\fn{true}$. وتمثل قيم الصدق
بدوال اختيار كما سبق؛ فإذا كان~$x$ قيمة صدق وطُبق في وسيطين، كانت
النتيجة الوسيط الأول إذا كان~$x$ هو~$\fn{true}$، والوسيط الثاني في
غير ذلك. وتأخذ~$\fn{And}$ وسيطيها~$x$ و$y$، ثم تمرر~$y$ و$\fn{false}$
إلى وسيطها الأول~$x$. وإذا افترضنا أن~$x$ قيمة صدق، قيّمت النتيجة
إلى~$y$ إذا كان~$x$ هو~$\fn{true}$، وإلى~$\fn{false}$ إذا كان~$x$
هو~$\fn{false}$، وهذا هو المطلوب بالضبط.

لاحظ أننا نفترض هنا ألا تستعمل مع~$\fn{And}$ إلا قيم الصدق. فإذا
مررت إليها حدود أخرى، فقد لا تكون النتيجة (أي الصورة السوية إن وجدت)
قيمة صدق.

\begin{prob}
  عرّف الدالتين~$\fn{Or}$ و$\fn{Xor}$ اللتين تمثلان دالتي صدق الفصل
  الشامل والفصل الحصري، مستعملًا ترميز قيم الصدق بحدود~$\lambd$.
\end{prob}

\end{document}
